\section{Methodology}

Our approach to studying $$-recursive sequence spaces combines techniques from functional analysis and computability theory. The methodology consists of two main components: theoretical analysis and computational validation.

\subsection{Theoretical Framework}

For a given recursive function $f: \mathbb{N} \to \mathbb{N}$, we analyze $_f$ through the following steps:

1. We first establish the topological properties of $_f$ by constructing a suitable metric:
\[
d(x,y) = \sum_{n=1}^{\infty} \frac{|x_n - y_n|}{2^n(1 + f(n))}
\]
where $x = (x_n)$ and $y = (y_n)$ are sequences in $_f$.

2. For each $_f$, we examine completeness by showing that Cauchy sequences converge within the space:
\[
\forall \varepsilon > 0, \exists N \in \mathbb{N}: m,n \geq N \implies d(x^{(m)}, x^{(n)}) < \varepsilon
\]

3. We investigate the algebraic structure by defining operations:
\[
(x + y)_n = x_n + y_n, \quad (\alpha x)_n = \alpha x_n
\]
and proving they preserve $f$-recursiveness.

\subsection{Computational Validation}

To verify our theoretical results, we implement the following:

1. A computer algebra system to generate explicit examples of $_f$ for various recursive functions $f$, including:
   - Polynomial functions
   - Exponential functions
   - Primitive recursive functions

2. Numerical experiments to test:
   - Convergence rates of sequences
   - Preservation of recursive bounds
   - Completeness properties

3. Implementation of algorithms to:
\[
\text{Verify}(x,f,n) = \begin{cases}
1 & \text{if } x_n \leq f(n) \\
0 & \text{otherwise}
\end{cases}
\]

All computational results are validated against theoretical predictions using rigorous error analysis with bounds:
\[
|E_n| \leq \frac{C}{f(n)} \quad \text{for some constant } C > 0
\]

This dual approach of theoretical analysis and computational validation ensures our results are both mathematically rigorous and practically verifiable.